\section{Discussion}
\label{sec:discussion}

\paragraph{What \sys{} measures.} A measurement is only as good as the link
between the construct and the task that operationalizes
it~\cite{jacobs2021measurement,raji2021benchmark}. \sys{} audits whether a
climate-policy study's identification assumptions are \emph{adequately supported
by its reported evidence}: reported-evidence adequacy is an operational proxy for
identification credibility, not causal validity itself. The two can diverge---a
study may report every diagnostic yet still be invalid, or omit a check yet be
sound---so \sys{} surfaces risk for an expert to weigh rather than adjudicating the
effect it estimates.

\paragraph{Evidence grounding is the binding constraint.} A further gap separates
the construct from the task: retrieving the relevant evidence. Our real-paper runs
show this, and not the reasoning, is the limiting factor. A reasoning assessor
handed off-target evidence over-flags every dimension, so \sys{} reports a
\emph{retrieval-quality} status and an explicit \texttt{unknown} risk (``the system
could not find the evidence'') distinct from a substantive high-risk judgement
(``the paper lacks it''). This is why the pipeline separates retrieval, relevance
filtering, and adequacy assessment.

\paragraph{Threats to validity.} These results are preliminary by design, spanning
synthetic fixtures, 27 real DID papers, and a human gold that is small (5 papers,
two annotators, one adjudication round); its two \emph{high} cells make high
precision/recall noisy, and the keyword baseline is an intentional lower bound.
Flaw injection certifies an auditor against the \emph{injected} threat on a single
dimension---it does not certify the absence of un-injected threats and inherits
whatever blind spots the flaw taxonomy has. We therefore treat detection and
localization as necessary but not sufficient evidence of an auditor's quality,
pairing them with the false-alarm rate and the human-gold comparison.

\paragraph{Future work.} Four directions follow from \S\ref{sec:results}:
\emph{(1)}~harden the reasoning assessor on subtler flaws, studying calibration,
cross-model robustness, and cited-evidence traceability under the same
flaw-injection protocol; \emph{(2)}~collect \emph{dimension-level} expert
annotation (the real corpus is labelled per method/claim, not per assumption) to
test whether \sys{}'s judgements match expert opinion; \emph{(3)}~run
human-expert studies of whether \sys{} reports actually help (trust, time,
errors avoided), since the system assists rather than replaces expert judgement;
and \emph{(4)}~extend beyond the China emissions-trading case to broader
environmental and empirical-economics settings, testing how far the
rubric-plus-injection methodology generalizes.
